\documentclass[11pt]{article}

\usepackage[T1]{fontenc}
\usepackage{lmodern}
\usepackage{amsmath,amssymb,amsthm}
\usepackage[margin=1in]{geometry}
\usepackage{microtype}
\usepackage{needspace}
\usepackage[hidelinks]{hyperref}
\hypersetup{
  pdftitle={A counterexample to Erdos Problem 1075},
  pdfauthor={Qiyuan Gu}
}

\newtheorem{theorem}{Theorem}
\newtheorem{proposition}[theorem]{Proposition}
\newtheorem{lemma}[theorem]{Lemma}
\numberwithin{equation}{section}
\newcommand{\W}{\mathcal W}
\newcommand{\D}{\mathcal D}

\title{A counterexample to Erd\H{o}s Problem 1075}
\author{Qiyuan Gu\\
University of Chicago\\
\href{mailto:phoenix1203@uchicago.edu}{\texttt{phoenix1203@uchicago.edu}}}
\date{}

\begin{document}
\maketitle

\begin{abstract}
Erd\H{o}s asked whether, for every fixed \(r\ge3\), there is a constant
\(c_r>r^{-r}\) such that, for every \(\varepsilon>0\), every
sufficiently large \(r\)-uniform hypergraph on \(N\) vertices with
at least \((1+\varepsilon)(N/r)^r\) edges contains
a subgraph on \(m\to\infty\) vertices with at least \(c_rm^r\) edges.
We give a counterexample for \(r=16\). We construct explicit finite
hypergraphs \(G_n\) whose unnormalized Lagrangians satisfy
\[
 16^{-16}\left(1+\frac1{4(14n)^{14}}\right)
 \le\lambda(G_n)\le16^{-16}+\frac1{15!n}.
\]
Blow-ups of these hypergraphs disprove the asserted density increase.
The construction uses two families of bipartite links arranged around
a cycle. A quantitative inequality for the corresponding open paths
gives the upper bound.
\end{abstract}

\section{Introduction}

All hypergraphs in this paper are finite and simple. For a hypergraph
\(H\), write \(v(H)=|V(H)|\), \(e(H)=|E(H)|\), and let \(H[S]\)
denote the subgraph induced by \(S\subseteq V(H)\).

Erd\H{o}s proved that, for every fixed \(r\ge2\) and \(\varepsilon>0\),
an \(r\)-uniform hypergraph on \(N\) vertices with at least
\(\varepsilon N^r\) edges contains a complete \(r\)-partite
\(r\)-uniform hypergraph with equal part sizes tending to infinity
with \(N\) \cite{Erdos64}. In particular, it contains a subgraph
on \(m\to\infty\) vertices with \((m/r)^r\) edges.
Erd\H{o}s subsequently conjectured that, for each fixed \(r\ge3\),
there is a constant \(c_r>r^{-r}\) such that, for every
\(\varepsilon>0\) and all sufficiently large \(N\), the assumption
\[
 e(H)\ge(1+\varepsilon)(N/r)^r
\]
forces a subgraph on \(m=m(N)\to\infty\) vertices with at least
\(c_rm^r\) edges \cite[p.~80]{Erdos74}.
This is Problem~1075 in Bloom's collection \cite{Bloom}.

\begin{theorem}\label{thm:main}
For every \(\gamma>16^{-16}\), there is an \(\varepsilon>0\) and
there are arbitrarily large \(16\)-uniform hypergraphs \(H\) such that
\[
 e(H)\ge(1+\varepsilon)\bigl(v(H)/16\bigr)^{16},
 \qquad
 e(H[S])<\gamma|S|^{16}
 \quad\text{for every nonempty }S\subseteq V(H).
\]
\end{theorem}

The theorem disproves the conjecture at the fixed uniformity \(r=16\).
It follows from a sequence of explicit finite hypergraphs with
Lagrangians approaching \(16^{-16}\) from above. We use the convention
\begin{equation}\label{eq:lagrangian}
 P_H(x)=\sum_{e\in E(H)}\prod_{v\in e}x_v,
 \qquad
 \lambda(H)=\max_{\substack{x_v\ge0\\\sum_vx_v=1}}P_H(x),
\end{equation}
without the factor \(r!\).

\begin{proposition}\label{prop:sequence}
For every integer \(n\ge2\), the \(16\)-uniform hypergraph \(G_n\)
defined in Section~\ref{sec:construction} has \(2n+29\) vertices and
satisfies
\begin{equation}\label{eq:bounds}
 16^{-16}\left(1+\frac1{4(14n)^{14}}\right)
 \le\lambda(G_n)\le16^{-16}+\frac1{15!n}.
\end{equation}
\end{proposition}

The construction couples two three-uniform hypergraphs by a cyclic
shift of their bipartite links. A prescribed weighting makes both
cubic forms attain their separate upper bounds, and one additional
edge gives a strict increase. When the additional edge has positive
weight, cutting the cycle prevents simultaneous equality in the two
cubic bounds. We quantify the resulting loss and use it to bound
every open path.
Deleting a pair of vertices of least total weight then gives the
upper bound in \eqref{eq:bounds}.

\section{The cyclic construction}\label{sec:construction}

Fix \(n\ge2\). Take pairwise distinct vertices
\[
 A_i,B_i\quad(i\in\mathbb Z/n\mathbb Z),\qquad C,U,V,
 \qquad W_1,\ldots,W_{12},\qquad D_1,\ldots,D_{14}.
\]
Put \(\W=\{W_1,\ldots,W_{12}\}\) and
\(\D=\{D_1,\ldots,D_{14}\}\).
On the vertex set \(\{A_i,B_i:i\in\mathbb Z/n\mathbb Z\}\cup
\{C\}\cup\D\), define two three-uniform hypergraphs by
\begin{align*}
 E(H_1)&=\bigcup_i
 \bigl\{\{A_i,X,Y\}:X\in\{B_i,C\},\
 Y\in\{B_j:j\ne i\}\cup\D\bigr\},\\
 E(H_2)&=\bigcup_i
 \bigl\{\{B_i,X,Y\}:X\in\{A_{i+1},C\},\
 Y\in\{A_j:j\ne i+1\}\cup\D\bigr\}.
\end{align*}
All subscripts are read modulo \(n\). Define
\begin{equation}\label{eq:edges}
 \begin{split}
 E(G_n)={}&\bigl\{\W\cup\{U\}\cup e:e\in E(H_1)\bigr\}\\
 &\cup\bigl\{\W\cup\{V\}\cup e:e\in E(H_2)\bigr\}
 \cup\bigl\{\{U,V\}\cup\D\bigr\}.
 \end{split}
\end{equation}
Each edge has sixteen vertices. The three displayed families are
disjoint, and the apex \(A_i\) or \(B_i\) determines the index within
each link family, so the construction has no repeated edges.

Use lower-case letters for the weights of the corresponding vertices,
and write
\[
 A=\sum_i a_i,\qquad B=\sum_i b_i,\qquad
 D=\sum_{j=1}^{14}d_j,\qquad \omega=\prod_{j=1}^{12}w_j.
\]
The two cubic forms and the polynomial of \(G_n\) are
\begin{align}
 P_1&=\sum_i a_i(b_i+c)(B-b_i+D),\label{eq:p1}\\
 P_2&=\sum_i b_i(a_{i+1}+c)(A-a_{i+1}+D),\label{eq:p2}\\
 P_{G_n}&=\omega(uP_1+vP_2)+uv\prod_{j=1}^{14}d_j.
 \label{eq:polynomial}
\end{align}

\Needspace{8\baselineskip}
\begin{lemma}\label{lem:lower}
The lower bound in \eqref{eq:bounds} holds.
\end{lemma}

\begin{proof}
Assign the weights
\begin{equation}\label{eq:weights}
 \begin{gathered}
 w_j=\frac1{16},\qquad u=v=\frac1{32},\qquad
 a_i=b_i=\frac1{16n},\\
 d_j=\frac1{224n},\qquad c=\frac{n-1}{16n}.
 \end{gathered}
\end{equation}
Their sum is one. Moreover,
\[
 A=B=c+D=\frac1{16},\qquad
 b_i+c=B-b_i+D=a_{i+1}+c=A-a_{i+1}+D=\frac1{16}.
\]
Thus \(P_1=P_2=16^{-3}\). The first term in
\eqref{eq:polynomial} is \(16^{-16}\), and the second is
\[
 \frac1{1024(224n)^{14}}
 =\frac{16^{-16}}{4(14n)^{14}}.\qedhere
\]
\end{proof}

\section{A bound for open paths}\label{sec:paths}

We first prove an inequality for the cubic forms obtained by opening
the cycle. The parameter \(D\) in the following lemma is a single
nonnegative number; its division among fourteen vertices will be
used only in Lemma~\ref{lem:path}.

\begin{lemma}\label{lem:cubic}
Let \(L\ge1\), and let \(a_1,\ldots,a_L,b_1,\ldots,b_L,c,D\)
be nonnegative numbers. Put
\[
 A=\sum_{i=1}^L a_i,\qquad B=\sum_{i=1}^L b_i,
 \qquad a_{L+1}=0,
\]
and suppose that \(A+B+c+D=1\). Define \(P_1,P_2\) by
\eqref{eq:p1} and \eqref{eq:p2}, with sums over \(1\le i\le L\).
Then, for every \(s\in[0,1]\),
\begin{equation}\label{eq:cubic-loss}
 sP_1+(1-s)P_2\le\frac1{27}-\frac{s(1-s)D^3}{1200}.
\end{equation}
\end{lemma}

\begin{proof}
Set
\[
 \delta_A=A-\frac13,\qquad\delta_B=B-\frac13,
\]
and
\begin{align*}
 \beta&=\frac{B+D-c}{2}=D+\frac{\delta_A}{2}+\delta_B,\\
 \alpha&=\frac{A+D-c}{2}=D+\delta_A+\frac{\delta_B}{2}.
\end{align*}
Completing the square gives
\begin{align}
 P_1&=\frac{A(1-A)^2}{4}-\sum_i a_i(b_i-\beta)^2,
 \label{eq:square1}\\
 P_2&=\frac{B(1-B)^2}{4}-\sum_i b_i(a_{i+1}-\alpha)^2.
 \label{eq:square2}
\end{align}
For \(0\le x\le1\),
\[
 \frac1{27}-\frac{x(1-x)^2}{4}
 =\frac{(x-1/3)^2(4/3-x)}4\ge\frac{(x-1/3)^2}{12}.
\]
Consequently, if
\[
 E=s\sum_i a_i(b_i-\beta)^2
 +(1-s)\sum_i b_i(a_{i+1}-\alpha)^2,
\]
then
\begin{equation}\label{eq:loss}
 \frac1{27}-sP_1-(1-s)P_2
 \ge\frac{s\delta_A^2+(1-s)\delta_B^2}{12}+E.
\end{equation}
This proves the assertion when \(D=0\) or \(s\in\{0,1\}\).
Assume henceforth that \(D>0\) and \(0<s<1\).

If \(|\delta_A|\ge D/10\), the right side of \eqref{eq:loss}
is at least
\[
 \frac{sD^2}{1200}\ge\frac{s(1-s)D^3}{1200},
\]
since \(D\le1\). The case \(|\delta_B|\ge D/10\) is analogous.
It remains to consider
\( |\delta_A|,|\delta_B|<D/10\).
In this case
\[
 D\le c+D=\frac13-\delta_A-\delta_B\le\frac13+\frac D5,
\]
so \(D\le5/12\). It follows that
\begin{equation}\label{eq:alpha-beta}
 \frac{17D}{20}\le\alpha,\beta\le\frac{23D}{20}\le\frac{23}{48},
 \qquad A,B\ge\frac13-\frac D{10}\ge\frac7{24}.
\end{equation}
In particular, \(\alpha,\beta>0\), \(A\ge\alpha/3\), and
\(B\ge\beta/3\).

Consider the finite sequence
\begin{equation}\label{eq:sequence}
 \frac{a_1}{\alpha},\frac{b_1}{\beta},\ldots,
 \frac{a_L}{\alpha},\frac{b_L}{\beta},0.
\end{equation}
Each successive pair \(z,z'\) contributes to \(E\) a term
\[
 s\alpha\beta^2 z(z'-1)^2
 \quad\text{or}\quad
 (1-s)\beta\alpha^2 z(z'-1)^2,
\]
according as the first entry comes from an \(a_i\) or a \(b_i\).
If some entry of \eqref{eq:sequence} is at least \(1/3\), then,
because the last entry is zero, some successive pair satisfies
\(z\ge1/3\) and \(z'<1/3\). Hence
\begin{equation}\label{eq:transition}
 E\ge\frac4{27}\min\{s\alpha\beta^2,(1-s)\beta\alpha^2\}.
\end{equation}
If all entries are less than \(1/3\), then instead
\begin{align*}
 E&\ge\frac49\bigl(s\beta^2A+(1-s)\alpha^2B\bigr)\\
 &\ge\frac4{27}\bigl(s\alpha\beta^2+(1-s)\beta\alpha^2\bigr),
\end{align*}
which also implies \eqref{eq:transition}. Using
\eqref{eq:alpha-beta}, we obtain in either case
\[
 E\ge\frac4{27}\left(\frac{17}{20}\right)^3
       \min\{s,1-s\}D^3
 \ge\frac{s(1-s)D^3}{1200}.
\]
The last inequality follows from
\(\frac4{27}(\frac{17}{20})^3=4913/54000>1/1200\).
Together with \eqref{eq:loss}, this proves the lemma.
\end{proof}

For \(L\ge1\), define the \(16\)-uniform hypergraph \(T_L\)
as in \eqref{eq:edges}, with each index ranging over \(1,\ldots,L\)
and with the last link in \(H_2\) replaced by the triples
\[
 \{B_L,C,Y\},\qquad Y\in\{A_1,\ldots,A_L\}\cup\D.
\]
The edge \(\{U,V\}\cup\D\) is retained.
Its polynomial is \eqref{eq:polynomial} with \(a_{L+1}=0\).

\begin{lemma}\label{lem:path}
For every \(L\ge1\), we have \(\lambda(T_L)\le16^{-16}\).
\end{lemma}

\begin{proof}
Take nonnegative vertex weights of total mass one, and put
\[
 h=\sum_{j=1}^{12}w_j,\qquad q=u+v,\qquad t=A+B+c+D.
\]
Thus \(h+q+t=1\). If \(q=0\) or \(t=0\), then \(P_{T_L}=0\).
Otherwise put
\[
 s=u/q,\qquad d=D/t,\qquad z=s(1-s)d^3\in[0,1/4].
\]
Applying Lemma~\ref{lem:cubic} to the bottom weights divided by
\(t\), and then using homogeneity, gives
\[
 uP_1+vP_2\le qt^3\left(\frac1{27}-\frac z{1200}\right).
\]
The arithmetic--geometric mean inequality now yields
\begin{equation}\label{eq:path-first}
 P_{T_L}\le
 \left(\frac h{12}\right)^{12}qt^3
       \left(\frac1{27}-\frac z{1200}\right)
 +q^2s(1-s)\left(\frac{td}{14}\right)^{14}.
\end{equation}
Also,
\begin{equation}\label{eq:amgm}
 \frac{qt^3}{27}\le\left(\frac{q+t}{4}\right)^4,
 \qquad
 \frac{q^2t^{14}}{14^{14}}\le4\left(\frac{q+t}{16}\right)^{16}.
\end{equation}
The second inequality applies the same inequality to two copies of
\(q/2\) and fourteen copies of \(t/14\).

Define
\[
 F(h)=\left(\frac{4h}{3}\right)^{12}[4(1-h)]^4.
\]
We have \(F(h)\le1\), with equality at \(h=3/4\).
Since \(1-27z/1200>0\), we may substitute \eqref{eq:amgm}
into \eqref{eq:path-first}. Using \(d^{14}\le d^3\), we get
\begin{align}
 \frac{P_{T_L}}{16^{-16}}
 &\le F(h)\left(1-\frac9{400}z\right)
        +4s(1-s)d^{14}(1-h)^{16}\notag\\
 &\le F(h)+z\left[4(1-h)^{16}-\frac9{400}F(h)\right].
 \label{eq:scalar}
\end{align}

Suppose first that \(1/2\le h<1\). Then
\[
 \frac{(1-h)^{16}}{F(h)}
 =\frac{(3/4)^{12}}{256}\left(\frac{1-h}{h}\right)^{12}
 \le\frac1{256}.
\]
Since \(1/64<9/400\), the expression in brackets in
\eqref{eq:scalar} is nonpositive. Thus the right side is at most
\(F(h)\le1\). If \(h=1\), the polynomial is zero.

Suppose next that \(0\le h\le1/2\). Discarding the negative term
in \eqref{eq:scalar} and using \(z\le1/4\), we obtain
\begin{equation}\label{eq:small-h}
 \frac{P_{T_L}}{16^{-16}}\le F(h)+(1-h)^{16}.
\end{equation}
For \(0<h\le1/2\), the function \(F(h)/h\) is increasing, since
its logarithmic derivative is
\[
 \frac{11}{h}-\frac4{1-h}=\frac{11-15h}{h(1-h)}>0.
\]
Moreover,
\[
 \frac{F(1/2)}{1/2}=32\left(\frac23\right)^{12}
 =\frac{131072}{531441}<1.
\]
It follows that \(F(h)\le h\) throughout \([0,1/2]\).
The right side of \eqref{eq:small-h} is therefore at most
\(h+(1-h)^{16}\le1\), as required.
\end{proof}

\section{The Lagrangian bound and the blow-ups}\label{sec:conclusion}

\begin{proof}[Proof of Proposition~\ref{prop:sequence}]
The lower bound was proved in Lemma~\ref{lem:lower}.
For the upper bound, take any nonnegative weights on \(G_n\) of
total mass one. Choose an index \(k\) with
\[
 a_k+b_k\le\frac{A+B}{n}\le\frac1n.
\]
Deleting \(A_k,B_k\) leaves an induced copy of \(T_{n-1}\): order
the remaining indices as \(k+1,\ldots,k+n-1\) modulo \(n\).
The last \(B\)-vertex loses its successor \(A_k\), giving precisely
the terminal link in the definition of \(T_{n-1}\).
By Lemma~\ref{lem:path} and homogeneity, the total weight of the
edges that remain is at most
\[
 16^{-16}(1-a_k-b_k)^{16}\le16^{-16}.
\]

For any simple \(16\)-uniform hypergraph and weights of total
mass one,
\begin{equation}\label{eq:derivative}
 \frac{\partial P_H}{\partial x_v}
 \le\sum_{\substack{S\subseteq V(H)\setminus\{v\}\\|S|=15}}
        \prod_{w\in S}x_w
 \le\frac{(\sum_{w\ne v}x_w)^{15}}{15!}
 \le\frac1{15!}.
\end{equation}
Indeed, each product of fifteen distinct weights occurs \(15!\)
times in the power expansion, and all other terms are nonnegative.
The deleted edges consequently have total weight at most
\[
 a_k\frac{\partial P_{G_n}}{\partial a_k}
 +b_k\frac{\partial P_{G_n}}{\partial b_k}
 \le\frac{a_k+b_k}{15!}\le\frac1{15!n}.
\]
Edges containing both deleted vertices are counted twice, which
does not affect the upper bound. This proves \eqref{eq:bounds}.
\end{proof}

\begin{proof}[Proof of Theorem~\ref{thm:main}]
Fix \(\gamma>16^{-16}\), and choose an integer \(n\ge2\) such that
\[
 16^{-16}+\frac1{15!n}<\gamma.
\]
Keep \(n\) fixed, and set
\[
 \varepsilon=\frac1{4(14n)^{14}}>0.
\]
\Needspace{9\baselineskip}
For each positive integer \(M\), form a blow-up \(B_{n,M}\) of
\(G_n\): replace each vertex by an independent class, and each
edge by all transversal edges across its sixteen classes.
Use the following class sizes, with one class for each listed vertex:
\[
 \begin{array}{c|ccccc}
 \text{vertex}&W_j&U,V&A_i,B_i&C&D_j\\ \hline
 \text{class size}&14nM&7nM&14M&14(n-1)M&M
 \end{array}.
\]
There are \(N=224nM\) vertices. These class sizes are proportional
to the weights in \eqref{eq:weights}, so Lemma~\ref{lem:lower}
gives the exact edge count
\begin{equation}\label{eq:blowup-count}
 e(B_{n,M})=\bigl[(14n)^{16}+49n^2\bigr]M^{16}
 =(1+\varepsilon)(N/16)^{16}.
\end{equation}

Let \(S\) be any nonempty set of \(m\) vertices of \(B_{n,M}\),
and let \(k_v\) be the number of its vertices in the class
corresponding to \(v\in V(G_n)\). Then
\[
 e(B_{n,M}[S])
 =\sum_{e\in E(G_n)}\prod_{v\in e}k_v
 =m^{16}P_{G_n}\bigl((k_v/m)_v\bigr)
 \le m^{16}\lambda(G_n)<\gamma m^{16}.
\]
Together with \eqref{eq:blowup-count} and \(M\to\infty\), this
proves the theorem. A non-induced subgraph has no more edges than
the induced subgraph on the same vertex set.
\end{proof}

\section*{AI tool disclosure}

OpenAI Codex (GPT-6) was used substantially in developing this work.
It proposed the cyclic construction, derived the quantitative path
inequality and the Lagrangian estimates, and drafted the manuscript.
It also assisted with literature searches and exact algebraic checks.
The author specified the research problem and directed the preparation
of the paper, and is responsible for its mathematical content.

\begin{thebibliography}{9}

\bibitem{Erdos64}
P.~Erd\H{o}s,
\href{https://users.renyi.hu/~p_erdos/1964-13.pdf}
{On extremal problems of graphs and generalized graphs},
\emph{Israel J. Math.} \textbf{2} (1964), 183--190.

\bibitem{Erdos74}
P.~Erd\H{o}s,
\href{https://users.renyi.hu/~p_erdos/1974-10.pdf}
{Extremal problems on graphs and hypergraphs},
in \emph{Hypergraph Seminar} (Columbus, Ohio, 1972),
Lecture Notes in Mathematics, vol.~411,
Springer, Berlin, 1974, pp.~75--84.

\bibitem{Bloom}
T.~F. Bloom,
\emph{Erd\H{o}s Problem \#1075},
\url{https://www.erdosproblems.com/1075}.
Accessed 5 September 2026.

\end{thebibliography}

\end{document}
